   (1): $m$ に関する数学的帰納法で「$2$でちょうど$m$回割れるようなすべての自然数$a$に対して$5^a - 1$が$2$でちょうど$m+2$回割れる」ことを証明する. $m=0$のとき, $a$は奇数であり, $5^a - 1 \equiv 4\pmod{8}$ であるから$2$でちょうど$2$回割れるから適する. $m=k\geq 0$ のときに主張が正しいと仮定して $m=k+1$の場合を示す. $a$は$2$でちょうど$k+1$回割れるとする. 偶数なので, $a = 2b$ とおけば
    $$
    5^a - 1 = (5^{b} - 1) (5^b + 1)
    $$
    であり, 仮定より$5^b- 1$は$2$でちょうど$k+2$回割れる. 一方$5^b + 1 \equiv 2\pmod{4}$より, $5^b+1$は$2$でちょうど$1$回割れるので, $5^a - 1$はちょうど$k+3 = (k+1) + 2$回割れる. よって$m=k+1$の場合も正しく, 題意は示された.


    (2): $\frac{10^K - N}{N^2}$が整数であるから, $10^K$ が $N$で割り切れる. よって$N=2^a5^b$ と表せるので,
    \[\frac{10^K - 2^a5^b}{2^{2a}5^{2b}}\]
    が整数となる($a,b\geq 0$).

    もし $a < K$であれば, 分子は$2$で$a$回割れるが, 分母は$2$で$2a$回割れるので, $a=0$でなければならない. 同様に$b < K$であれば$b=0$である.

    また $a > K$であれば, $10^K - 2^a5^b$は$2$で$K$回割り切れるが, $2^{2a}$は$2$で$2K$回以上割り切れるので不適である. 同様に $b>K$ も不適である.

    よって, $a,b$は$0$または$K$である.

    $a=b=0$のとき$N=1$なので$N\geq 2$に反する.

    $(a,b) = (K,K)$のときは$N=10^K$で, $\frac{10^K - 10^K}{10^{2K}} = 0$ なので, 組 $(N,K) = (10^i, i)$ ($i=1,2\dots$) は適する.

    $(a,b) = (K,0)$のとき, $N = 2^K$であり, $\frac{10^K - 2^K}{2^{2K}} = \frac{5^K - 1}{2^K}$なので, $5^{K}-1$が$2$で$K$回割り切れる. よって(1)より, $K$が$2$で$K-2$回割れる. よって $K \geq 2^{K-2}$であることが必要なので, $K=1,2,3,4$ が必要であり(証明略), 実際に計算すると
    \[
    \frac{5^K - 1}{2^K} = 2, 6, \frac{31}{2}, 39
    \]
    なので$K=1,2,4$が適する. よって組 $(N,K) = (2,1), (4,2), (16,4)$を得る.

    $(a,b) = (0,K)$のとき, $n=5^K$であり, $\frac{10^K - 5^K}{5^{2K}} = \frac{2^K - 1}{5^K}$となるが, $5^K > 2^K - 1 > 0$なので整数になることはない.

    以上より, 求める$(N,K)$は
    \[
    (N,K) = (2,1), (4,2), (16,4), (10^i, i) \quad (i=1,2,\dots)
    \]   
